\documentclass[11pt]{article}
\usepackage[margin=1in]{geometry}
\usepackage{amsmath,amssymb,amsthm}
\usepackage{hyperref}
\usepackage{microtype}

\newtheorem{theorem}{Theorem}
\newtheorem{lemma}[theorem]{Lemma}

\title{Counterexamples to a Source-Pair Augmentation Question\\for Boolean and Binary Rank}
\author{Keigo Oka}
\date{August 30, 2026}

\begin{document}
\maketitle

\begin{abstract}
Parnas and Shraibman asked whether, when two source bases witness failure of the augmentation property for Boolean or binary rank, one can always choose one vector from each source base so that augmenting by only those two vectors already increases the rank. We answer the question negatively for both ranks. The Boolean counterexample has four rows; the binary counterexample has five rows and is row-minimal among counterexamples of this form. The core proofs are elementary. An accompanying exhaustive verifier independently checks the base graphs, all cross-pair augmentations, and the binary row-minimality statement.
\end{abstract}

\section{Introduction}

For a binary matrix, Parnas and Shraibman introduced the augmentation property for Boolean and binary rank and characterized it in terms of the directed graph of optimal bases \cite{PS}. If a matrix has at least two source bases, augmenting by all vectors from two sources raises the rank. In Section~6 they ask whether this failure is always witnessed more locally: given two source bases $U,V$, must there exist $u\in U$ and $v\in V$ such that adjoining only $u,v$ already raises the rank?

We show that the answer is no for both ranks.

We encode a binary column by its bitmask, with the least significant bit corresponding to row~1. For Boolean rank, spanning means closure under coordinatewise OR. For binary rank, spanning is by ordinary $0/1$ linear combinations; any selected summands in a representation of a binary target necessarily have pairwise disjoint supports.

A base is a minimum-cardinality spanning set. The base graph has one vertex for each base and an arrow $B\to C$ when $B$ spans every vector of $C$; a source has no incoming arrow from a distinct base.

\section{Boolean rank}

Let
\[
A=\{3,7,15\},\qquad U=\{3,5,8\},\qquad V=\{3,5,12\}.
\]

\begin{theorem}
For Boolean rank,
\[
\operatorname{rank}_{\mathrm B}(A)=3,
\qquad
\operatorname{rank}_{\mathrm B}(A\mid U\mid V)=4,
\]
$U,V$ are distinct source bases of $A$, and
\[
\operatorname{rank}_{\mathrm B}(A\mid u\mid v)=3
\qquad\text{for every }u\in U,\ v\in V.
\]
\end{theorem}

\begin{proof}
The sets $U$ and $V$ span $A$ because
\[
7=3\vee5,
\qquad
15=3\vee5\vee8=3\vee5\vee12.
\]
The columns of $A$ give rank at most three. Two Boolean generators have at most the three nonzero OR-values $p,q,p\vee q$. If $p,q$ are incomparable these values are not a chain; if they are comparable there are at most two distinct nonzero values. Hence two generators cannot span the strict chain
\[
3<7<15,
\]
so the rank is three.

To span $U=\{3,5,8\}$ with three Boolean generators, the singleton fourth-row vector $8$ must itself be a generator. It cannot be used in an OR equal to $3$ or $5$, so the other two generators must be $3$ and $5$. Thus $U$ is the unique three-generator spanning set of itself.

For $V=\{3,5,12\}$, producing $3$ requires a generator carrying row~2 and no rows outside $\{1,2\}$, while producing $12$ requires a distinct generator carrying row~4 and no rows outside $\{3,4\}$. Neither can be used in an OR equal to $5=(1,0,1,0)^{\mathsf T}$, so the third generator is $5$. The first two are then forced to be $3$ and $12$. Thus $V$ is also the unique three-generator spanning set of itself. In particular both are source bases of $A$.

If $A\cup U\cup V$ had Boolean rank three, a Boolean base of $A$ would span $U$, hence would equal $U$; but $U$ does not span $12$. Thus the rank is at least four, while the four coordinate vectors give an upper bound of four.

For the cross pairs, if $v\in\{3,5\}$ then $U$ spans $A\cup\{u,v\}$ for every $u\in U$. If $u\in\{3,5\}$ and $v=12$, then $V$ spans the augmentation. The remaining pair $(8,12)$ is spanned together with $A$ by
\[
W=\{3,4,8\},
\]
since $7=3\vee4$, $12=4\vee8$, and $15=3\vee4\vee8$. Every cross-pair augmentation therefore has rank three.
\end{proof}

\section{Binary rank}

Let
\[
A=\{10,31,27,18\},\qquad
U=\{4,9,10,18\},\qquad
V=\{9,10,18,21\}.
\]

\begin{theorem}
For binary rank,
\[
\operatorname{rank}_{\mathrm{bin}}(A)=4,
\qquad
\operatorname{rank}_{\mathrm{bin}}(A\mid U\mid V)=5,
\]
$U,V$ are distinct source bases of $A$, and
\[
\operatorname{rank}_{\mathrm{bin}}(A\mid u\mid v)=4
\qquad\text{for every }u\in U,\ v\in V.
\]
\end{theorem}

\begin{proof}
The real rank of the five-by-four matrix with columns $A$ is four: the minor on the first four rows has determinant $-1$. Hence binary rank is at least four, and the four displayed columns give equality.

The base $U$ spans $A$ via the support-disjoint sums
\[
27=9+18,
\qquad
31=4+9+18,
\]
while $V$ spans $A$ via
\[
27=9+18,
\qquad
31=10+21.
\]

We show first that $U$ is the unique four-generator binary spanning set of itself. In support notation,
\[
4=\{3\},\quad
9=\{1,4\},\quad
10=\{2,4\},\quad
18=\{2,5\}.
\]
To form the singleton $\{3\}$, one generator is forced to be exactly $\{3\}$. The other three generators must span $\{1,4\}$, $\{2,4\}$, and $\{2,5\}$. The first and third targets are disjoint. If either were split into two nonzero generators, their two decompositions together would consume all three remaining generators. The middle target intersects both and would have to reuse a generator from each side; but a one-generator side would import the unwanted coordinate $1$ or $5$, while splitting both disjoint targets would require at least four generators. Hence the two disjoint targets are themselves generators, and the remaining generator is exactly $\{2,4\}$. Thus the spanning set is $U$.

Next consider $V$. The same disjoint-target argument shows that any three generators spanning
\[
9=\{1,4\},\qquad10=\{2,4\},\qquad18=\{2,5\}
\]
must be exactly $9,10,18$. The target
\[
21=\{1,3,5\}
\]
contains row~3, whereas the first three do not. Any generator containing row~3 cannot be used for those three targets, so they already force three row-3-free generators $9,10,18$. None of them can participate in a support-disjoint representation of $21$ because each carries an unwanted coordinate. The fourth generator is therefore $21$. Thus $V$ is also unique as a four-generator spanning set of itself, and both $U,V$ are source bases.

If $A\cup U\cup V$ had binary rank four, a four-element base spanning it would span $U$, hence would equal $U$; but $U$ does not span $21$. Therefore the rank is at least five. The five coordinate unit vectors give rank at most five.

Finally,
\[
U\cap V=\{9,10,18\}.
\]
If $v$ lies in this intersection, $U$ spans $A\cup\{u,v\}$; if $u$ lies in it, $V$ spans the augmentation. The only remaining pair is $(u,v)=(4,21)$. It is spanned together with $A$ by
\[
W=\{4,10,17,18\},
\]
because
\[
21=4+17,
\qquad
27=10+17,
\qquad
31=4+10+17
\]
are support-disjoint sums. Hence all sixteen cross-pair augmentations have binary rank four.
\end{proof}

\section{Exact verification and row minimality}

The accompanying program \texttt{verification/verify\_source\_pair\_augmentation.py} enumerates all optimal bases and source bases for the two displayed matrices and checks all cross pairs using integer arithmetic only. It additionally exhausts all binary column sets on at most four rows and finds no binary-rank counterexample of this source-pair form. Thus the five-row binary example is row-minimal.

The computations are not needed for the two counterexample theorems themselves; they are included as an independent certificate and for the minimality statement.

\section{Discovery provenance}

The examples were found while experimenting with the Interaction Reconstruction programme. The broader programme is not used in either proof. The relevant heuristic was to regard the base graph as a finite factorization atlas and search for source-interface incompatibility that is jointly expensive but pairwise invisible.

\begin{thebibliography}{9}
\bibitem{PS}
M.~Parnas and A.~Shraibman,
\newblock The augmentation property of binary matrices for the binary and Boolean rank,
\newblock \emph{Linear Algebra and its Applications} 556 (2018), 70--99.
\newblock doi:10.1016/j.laa.2018.07.001.
\end{thebibliography}

\end{document}
